% !TeX root = ../../Book1.tex
% 纲要标题：### 8.4* 自由群的进一步性质
% 纲要位置：计划纲要.md，第 348 行。

\OptionalSection{自由群的进一步性质}{b1:sec:0804}

% 纲要内容（仅作写作计划，尚非教材正文）：
%
% * Nielsen–Schreier 定理的图论证明路线
% * 子群秩与 Schreier 公式的有限指数情形
% * 字问题与几何群论的入口
%
% ---
%

\subsection{Nielsen–Schreier 定理的图论证明}
\label{b1:subsec:0804:01}

本节给出\term[nielsen-schreier-dingli]{Nielsen--Schreier 定理}[Nielsen--Schreier theorem]的组合图论证明。所需的路径只是一串相接的有向边，不需要拓扑空间或基本群的先修知识。图论方法还会实际产生子群的自由基。

\ProseParagraphBreak
先约定图允许重边和环。每条几何边有两个相反方向 \(e,\bar e\)，即使是环也区分方向。一条边路径是首尾相接的有向边序列；插入或删除 \(e\bar e\) 称为消去一次往返。固定顶点 \(o\)，从 \(o\) 回到 \(o\) 的路径在往返消去的等价关系下，以串接为乘法组成群，记为 \(L(\Gamma,o)\)。结合律由串接给出，逆是倒序反向的路径。

\begin{lemma}\label{b1:lem:graph-loop-free}
设 \(\Gamma\) 是连通图，\(T\) 是包含全部顶点的树。为每条不在 \(T\) 中的几何边选一个方向，组成集合 \(E_0\)。从 \(o\) 到顶点 \(v\) 的唯一树路径记为 \(p_v\)。则 \(L(\Gamma,o)\) 是自由群，自由基为
\[
\ell_e=p_v\,e\,p_w^{-1}\qquad
(e\in E_0,\ e:v\longrightarrow w).
\]
\end{lemma}
\begin{proof}
树的两顶点间有唯一路径：存在性来自连通，若有两条不同的无往返路径，取首次分离与再次汇合的区段便得到回路，违背树的无回路性。树内的任何闭路径均可消为常路径；否则有限条边的路径中保留一条无往返闭路径，便给出回路。

把 \(E_0\) 当作自由字母，由泛性质得到 \(F(E_0)\to L(\Gamma,o)\)，字母 \(e\) 映为 \(\ell_e\)。反向定义一个映射：沿闭路径读取边，树边略过，非树边按所选方向读作 \(e\)，反向读作 \(e^{-1}\)，最后约化。插入往返时，若是树边不会读出字母，若是非树边则增加相邻逆字母，因此映射良定义且为同态。

在一个基环路 \(\ell_e\) 上，反向读取只得到字母 \(e\)，故一个方向的复合为恒等。为检查另一方向，设闭路径依次走边 \(e_i:v_{i-1}\to v_i\)，其中 \(v_0=v_k=o\)。在每个相邻位置插入 \(p_{v_i}^{-1}p_{v_i}\)，得到它等价于
\[
\prod_{i=1}^k
\bigl(p_{v_{i-1}}e_i p_{v_i}^{-1}\bigr).
\]
树边对应的括号项可在树中消去；非树边对应基环路或其逆。因此读出字母再送回环路，恰好恢复原来的等价类。两个同态互逆。
\end{proof}
连通图总可选一棵这样的生成树。有限图可以从根开始，每次添加通向新顶点的边。一般无限图在通常的选择公理下，对包含全部顶点的无回路子图用 Zorn 引理取极大者：链的并仍无回路，因为回路只用有限条边；极大者若不连通，可沿原图连接两个分支的路径加进一条跨分支边而不产生回路，矛盾。

\begin{theorem}[Nielsen--Schreier]\label{b1:thm:nielsen-schreier}
自由群的每个子群都是自由群。
\end{theorem}
\begin{proof}
设 \(H\le F(S)\)。构造以右陪集 \(Hg\) 为顶点的\term[schreier-peiji-tu]{Schreier 陪集图}[Schreier coset graph]：对每个顶点 \(Hg\) 和 \(s\in S\)，画一条指定正向为
\[
Hg\xrightarrow{s}Hgs
\]
的几何边，其反向标记为 \(s^{-1}\)。右乘使边与代表元选择无关。即使两条边有同样的端点，也保留其不同身份。任意字从顶点 \(H\) 出发有唯一的带相应字母标签的路径，其终点为该字代表的陪集，所以图连通。

以 \(H\) 为根。闭路径的标签乘积属于 \(H\)，并给出同态 \(L(\Gamma,H)\to H\)。每个 \(h\in H\) 的字都有回到根的路径，故此同态满射。若闭路径的标签在自由群中为单位，则其字可以相邻逆字母消为空字。图中每个顶点都有唯一的给定字母出边，所以路径上相邻逆字母必是同一条边的正反向。字的每次消去因此对应路径的一次往返消去，最终路径也消为空路径，同态单射。

选生成树并应用前一引理，得到 \(H\cong L(\Gamma,H)\cong F(E_0)\)。若树路径 \(p_v\) 的标签字为 \(t_v\)，基元素可显式写为
\[
t_v s\,t_w^{-1}
\]
其中 \(v\xrightarrow{s}w\) 取遍不在树中的正向边。这些字属于 \(H\)，因为两条到达 \(w\) 的路径 \(t_vs,t_w\) 表示相同右陪集。
\end{proof}
证明同时解释了自由群子群为何不会凭空产生新关系：子群中的字关系，转化为图中闭路径的往返消去；树消除了连接不同顶点的冗余路径，剩下的每条边恰好提供一个自由字母。

\subsection{有限指数子群与 Schreier 公式}
\label{b1:subsec:0804:02}

\term[schreier-zhi-gongshi]{Schreier 秩公式}[Schreier rank formula]把有限指数子群的自由基个数化为图的边数与顶点数之差。

\begin{theorem}[Schreier 秩公式]\label{b1:thm:schreier-rank}
设 \(F_r\) 的秩 \(r\) 有限，且子群 \(H\le F_r\) 的指数为有限数 \(d\)。则
\[
\operatorname{rank}H=1+d(r-1).
\]
\end{theorem}
\begin{proof}
Schreier 图有 \(d\) 个顶点，每个顶点对于每个基字母恰有一条正向边，故共有 \(dr\) 条几何边；反向边不重复计数。一棵含全部顶点的树有 \(d-1\) 条边。这一计数可对有限树逐次删去末端顶点归纳得到。非树边数于是为 \(dr-(d-1)\)，而它就是前面所构造自由基的大小。\(r=0\) 时只能有 \(d=1\)，公式仍成立。
\end{proof}
例如 \(F_2=F(\{a,b\})\) 到加法群 \(C_2\) 的同态 \(a\mapsto1,b\mapsto0\)，核 \(H\) 的指数为二。图有顶点 \(H,Ha\)，选从 \(H\) 到 \(Ha\) 的那条正向 \(a\) 边为树。余下正向边分别是 \(Ha\) 到 \(H\) 的另一条 \(a\) 边，以及两个顶点上的 \(b\) 环。树代表字为 \(1,a\)，所以所得自由基为
\[
a^2,\qquad b,\qquad aba^{-1}.
\]
公式也给出 \(\operatorname{rank}H=3\)。它表明非交换自由群子群的秩可以大于母群的秩，不能用向量空间子空间的维数直觉判断。

\subsection{字问题与几何群论}
\label{b1:subsec:0804:03}

给定一个有限生成群的呈示，\term[zi-wenti]{字问题}[word problem]要求判定输入字是否表示单位。对自由群，约化算法读入每个字母时至多进栈或出栈一次，因而在通常按字母计费的模型下，可用与字长成正比的步骤判定。对前面二面体群及四面体群的具体呈示，也可以按已证明的标准形进行判断。

\ProseParagraphBreak
一般呈示的关系闭包没有约化字那样的唯一代表。把关系字及其共轭反复插入、删除，可以验证某些相等式；若无法找到变换过程，不能据此断定两个字不相等，更不能保证搜索总会停止。本节的算法结论只针对已经给出标准形及终止论证的情形，不把它推广到任意有限呈示。

\ProseParagraphBreak
取 \(H=1\) 的 Schreier 图称为这里的\term[cayley-tu]{Cayley 图}[Cayley graph]：顶点为群元素，边记录右乘生成元。自由群的 Cayley 图是一棵树，因为无往返闭路径的标签会是非空约化字，不可能表示单位。给每条边长度一以后，从单位到一个元素的距离恰为其约化字长。于是代数乘法给出的字与图上的路径长度发生联系，这构成几何群论研究的一个基本出发点。更多关于自由群、呈示与组合方法的材料可参阅 Lyndon 与 Schupp 的《Combinatorial Group Theory》\cite{LyndonSchupp2001}。
